\documentclass[11pt]{article}
\usepackage[a4paper,margin=18mm]{geometry}
\usepackage{fontspec}
\setmainfont{Latin Modern Roman}
\usepackage{amsmath,amssymb,amsthm,amscd,graphicx,color}
\usepackage{xurl}
\usepackage[unicode,hidelinks]{hyperref}
\allowdisplaybreaks
\setlength\emergencystretch{3em}

\makeatletter
\newtheorem{theorem}{Theorem}
\newtheorem*{theorem*}{Theorem}
\newtheorem{lemma}{Lemma}
\newtheorem*{lemma*}{Lemma}
\newtheorem{corollary}{Corollary}
\newtheorem*{corollary*}{Corollary}
\newtheorem{proposition}{Proposition}
\newtheorem*{proposition*}{Proposition}
\newtheorem{conjecture}{Conjecture}
\newtheorem*{conjecture*}{Conjecture}
{\theoremstyle{definition} \newtheorem{definition}{Definition}
\newtheorem*{definition*}{Definition}
\newtheorem{example}{Example}
\newtheorem*{example*}{Example}
\newtheorem{remark}{Remark}
\newtheorem*{remark*}{Remark}
\newtheorem{note}{Note}
\newtheorem*{note*}{Note}
}
\makeatother
\usepackage{euscript}

%For tables:
\newcommand{\head}[1]{\textbf{#1}}

\newcommand{\McKay}{M$^{\text{c}}$Kay\xspace}
\newlength{\setBracketHeight}
\newcommand{\SetSuchThat}[2]{
  \settoheight{\setBracketHeight}{\ensuremath{#2}}
  \ensuremath{\left\{\left.{#1\rule{0cm}{\setBracketHeight}}\,
      \right|\,{#2}\right\}}}
\newcommand{\Span}[1]{\ensuremath{ \operatorname{span} \left \{ {#1}
    \right \} }}
% Lie derivative:
\newcommand{\LieDer}{\ensuremath{\EuScript L}}
% Hooking a vector into a differential form:
\newcommand{\hook}{\ensuremath{\mathbin{ \hbox{\vrule height1.4pt
        width4pt depth-1pt \vrule height4pt width0.4pt depth-1pt}}}}
\newcommand{\openSet}{\textrm{ open }}
\newcommand{\pd}[2]{\ensuremath{\frac{\partial{#1}}{\partial{#2}}}}
\newcommand{\R}[1]{\ensuremath{\mathbb{R}^{#1}}}
\newcommand{\C}[1]{\ensuremath{\mathbb{C}^{#1}}}
\newcommand{\Q}[1]{\ensuremath{\mathbb{Q}^{#1}}}
\newcommand{\Z}[1]{\ensuremath{\mathbb{Z}^{#1}}}
\newcommand{\Ha}[1]{\ensuremath{\mathbb{H}^{#1}}}
\newcommand{\Oc}[1]{\ensuremath{\mathbb{O}^{#1}}}
\newcommand{\CP}[1]{\ensuremath{\mathbb{CP}^{#1}}}
\newcommand{\RP}[1]{\ensuremath{\mathbb{RP}^{#1}}}
\newcommand{\HP}[1]{\ensuremath{\mathbb{HP}^{#1}}}
\newcommand{\OP}[1]{\ensuremath{\mathbb{OP}^{#1}}}
\newcommand{\Gr}[2]{\ensuremath{\operatorname{Gr}_{#1}\mathbb{C}^{#2}}}
\newcommand{\Fl}[2]{\ensuremath{\mathbb{F}_{#1}\mathbb{C}^{#2}}}
\newcommand{\LagFl}[2]{\ensuremath{\operatorname{Lag}_{#1}\mathbb{C}^{#2}}}
\newcommand{\NullFl}[2]{\ensuremath{\operatorname{Null}_{#1}\mathbb{C}^{#2}}}
%\newcommand{\Gr}[2]{\ensuremath{\operatorname{Gr}\left({#1},{#2}\right)}}
\newcommand{\Gro}[2]{\widetilde{\operatorname{Gr}}\left({#1},{#2}\right)}
\newcommand{\GrC}[2]{\operatorname{Gr}_{\C{}}\left({#1},{#2}\right)}
\newcommand{\Grnull}[2]{
\ensuremath{
\operatorname{Gr}_{\operatorname{null}}
\left({#1},{#2}\right)
}
}
\newcommand{\SO}[1]{\operatorname{SO}\left({#1}\right)}
\newcommand{\Spin}[1]{\ensuremath{\operatorname{Spin}\left({#1}\right)}}
\newcommand{\Cl}[1]{\ensuremath{\operatorname{C\ell}\left({#1}\right)}}
\newcommand{\PO}[1]{\operatorname{\mathbb{P}O}\left({#1}\right)}
\newcommand{\GL}[1]{\operatorname{GL}\left({#1}\right)}
\newcommand{\GLp}[1]{\operatorname{GL}^+\left({#1}\right)}
\newcommand{\SL}[1]{\operatorname{SL}\left({#1}\right)}
\newcommand{\PSL}[1]{\mathbb{P}\SL{#1}}
\newcommand{\PGL}[1]{\mathbb{P}\GL{#1}}
\newcommand{\Un}[1]{\operatorname{U}\left({#1}\right)}
\newcommand{\SU}[1]{\operatorname{SU}\left({#1}\right)}
\newcommand{\Symp}[1]{\ensuremath{\operatorname{Sp}\left({#1}\right)}}
\newcommand{\Sym}[2]{\ensuremath{\operatorname{Sym}^{#1}\left({#2}\right)}}
\newcommand{\Lm}[2]{\ensuremath{\Lambda^{#1} \left ( {#2} \right )}}
\newcommand{\nForms}[2]{\ensuremath{\Omega^{#1} \left ( {#2} \right
    )}}
\newcommand{\Cohom}[2]{\ensuremath{ H^{#1} \left ( {#2} \right
    )}}
\newcommand{\Homol}[2]{\ensuremath{ H_{#1} \left ( {#2} \right
    )}}
\DeclareMathOperator{\Ad}{Ad}
% \newcommand{\ModSp}[2]{\mathfrak{M}_{#1} \left ( {#2} \right )}
\newcommand{\ModSp}[2]{\ensuremath{{#2}^{#1}}}
\newcommand{\trans}[1]{{}^t{#1}} \newcommand{\bari}{\bar{\imath}}
\newcommand{\barj}{\bar{\jmath}}
\newcommand{\LT}[1]{\Lambda^3_+\left({#1}\right)}
\DeclareMathOperator{\Det}{Det}
\DeclareMathOperator{\tr}{tr}
\renewcommand{\Re}{\operatorname{Re}}
\renewcommand{\Im}{\operatorname{Im}}
\DeclareMathOperator{\ad}{ad}
\DeclareMathOperator{\Aut}{Aut}
\newcommand{\Lag}[1]{\ensuremath{\operatorname{Lag}\left(#1\right)}}
\newcommand{\Norm}[2]{\nu_{#2}{#1}}
\newcommand{\Conorm}[2]{\nu^*_{#2}{#1}}
\newcommand{\Proj}[1]{\mathbb{P}^{#1}}
\newcommand{\ucProj}[1]{\widetilde{\mathbb{P}}^{#1}}
\newcommand{\Aff}[1]{\mathbb{A}^{#1}}
\DeclareMathOperator{\ind}{ind}
\newcommand{\OO}[1]{
  \ensuremath{
    \mathcal{O}
    \ifthenelse{\equal{#1}{0}}
      {}
      {\left({#1}\right)}
  }
}
\newcommand{\OOp}[2]{
  \ensuremath{
    \mathcal{O}
    \ifthenelse{\equal{#1}{0}}
      {}
      {\left({#1}\right)}
    \ifthenelse{\equal{#2}{1}}
      {}
      {^{\oplus{#2}}}
  }
}
%\newcommand{\Bl}[2]{ % Blowup
%  \ensuremath{
%    \operatorname{Bl}_{#1}
%    \left(
%    {#2}
%    \right)
%  }
%}

\newcommand{\extraSection}[2]
{
\ifthenelse{\boolean{abridged}}
  {
  }
  {
    \section{#1}
    \begin{center}
    \emph{This section will not be referred to
    subsequently, and may be skipped.}
    \end{center}
    \par{#2}
  }
}
\newcommand{\extraSubsection}[2]
{
\ifthenelse{\boolean{abridged}}
  {
  }
  {
    \subsection{#1}
    \begin{center}
    \emph{This subsection will not be referred to
    subsequently, and may be skipped.}
    \end{center}
    \par{#2}
  }
}
\newcommand{\extraStuff}[1]
{
\ifthenelse{\boolean{abridged}}
  {
  }
  {
    {#1}
  }
}
\def\cprime{$'$}

%% Commands special to this document
\newcommand{\LieG}{\ensuremath{\mathfrak{g}}}
\newcommand{\LieP}{\ensuremath{\mathfrak{p}}}
\newcommand{\LieM}{\ensuremath{\mathfrak{m}}}
\newcommand{\LieA}{\ensuremath{\mathfrak{a}}}
\newcommand{\LieN}{\ensuremath{\mathfrak{n}}}
\newcommand{\LieH}{\ensuremath{\mathfrak{h}}}
\newcommand{\LieL}{\ensuremath{\mathfrak{l}}}
\newcommand{\sa}[1]{\ensuremath{\mathfrak{s}^{#1}}}
\newcommand{\SA}[1]{\ensuremath{S^{#1}}}
\newcommand{\Or}[1]{\ensuremath{\operatorname{Or}\left(#1\right)}}
\newcommand{\Gtot}{\ensuremath{G}}
\newcommand{\gtot}{\mathfrak{g}}
\newcommand{\Gstruc}{P}
\newcommand{\gstruc}{\mathfrak{p}}
\newcommand{\SSPart}{L}
\newcommand{\sspart}{\mathfrak{l}}
\newcommand{\AbPart}{A}
\newcommand{\abpart}{\mathfrak{a}}
\newcommand{\RedPart}{LA}
\newcommand{\redpart}{\mathfrak{l} \oplus \mathfrak{a}}
\newcommand{\NilPart}{N}
\newcommand{\nilpart}{\ensuremath{\mathfrak{n}}}
\newcommand{\nil}[1]{\ensuremath{{#1}^{\NilPart}}}
\newcommand{\semib}[1]{\ensuremath{{#1}^{\Gtot/\Gstruc}}}
\newcommand{\Cartan}{\ensuremath{H}}
\newcommand{\cartan}{\ensuremath{\mathfrak{h}}}
\newcommand{\Borel}{\ensuremath{B}}
\newcommand{\borel}{\ensuremath{\mathfrak{b}}}
\newcommand{\cpt}[1]{%
        \ifthenelse{\equal{#1}{\RedPart} \or \equal{#1}{\redpart}}
        {\ensuremath{(#1)^c}}
        {\ensuremath{#1^c}}
}
\newcommand{\red}[1]{\ensuremath{{#1}^{\RedPart}}}
\newcommand{\Flag}[1]{\ensuremath{\mathbb{F}\left({#1}\right)}}
\newcommand{\TT}{{}^{\sharp} t}
\newcommand{\NA}{\ast}
\newcommand{\Dual}{\Lambda}
\newcommand{\Hom}[2]{\ensuremath{\operatorname{Hom}\left(#1,#2\right)}}
\newcommand{\End}[1]{\ensuremath{\operatorname{End}\left(#1\right)}}
\newcommand{\rivf}[1]{\overleftarrow{#1}}
\newcommand{\livf}[1]{\overrightarrow{#1}}
\newcommand{\kernel}{\mathfrak{k}}
\newcommand{\Kernel}{K}
\newcommand{\s}[1]{\slLie{2,\C{}}_{#1}}
\newcommand{\Sl}[1]{\SL{2,\C{}}_{#1}}
\newcommand{\Bl}[1]{B_{#1}}
\newcommand{\No}[1]{N_{#1}}
\newcommand{\SpinorPlus}[1]{
\ensuremath{
\mathbb{S}^{+}\left({#1}\right)
}
}
\newcommand{\SpinorMinus}[1]{
\ensuremath{
\mathbb{S}^{-}\left({#1}\right)
}
}
\newcommand{\OPC}
{
\ensuremath{
\mathbb{OP}^2\left(\C{}\right)
}
}
\newcommand{\OPCdual}
{
\OPC^{\times}
}
\newcommand{\Grw}{
        \ensuremath{
        \operatorname{Gr}_{\Omega}\left(\mathbb{O}^3,\mathbb{O}^6\right)
        }
}
\DeclareMathOperator{\Tor}{Tor}
\DeclareMathOperator{\GW}{GW}

%% Associated graded
\DeclareMathOperator{\gr}{gr}

% Weight lattices:
\newcommand{\regWt}{}
\newcommand{\spclWt}{}

% Parabolic subgroups
\newcommand{\pin}{\bullet}
\newcommand{\pout}{\times}
\usepackage{ifthen,xspace}

\begin{document}
\begin{center}\Large Every holomorphic parabolic geometry with effective adjoint model on a compact complex manifold with trivial holomorphic tangent bundle descends from a left invariant Cartan geometry on a complex Lie group\end{center}
\noindent\textbf{Stable ID:} 2648\par
\noindent\textbf{Authors:} Benjamin McKay\par
\noindent\textbf{Original work:} Holomorphic Parabolic Geometries and Calabi-Yau Manifolds\par
\noindent\textbf{Version:} Official SIGMA 7 (2011) article 090, DOI 10.3842/SIGMA.2011.090; journal PDF and native TeX acquired 2026-10-01, exact bytes pinned by SHA256\par
\noindent\textbf{Selected task scope:} Theorem3 restricted to effective rational homogeneous models G/P with connected complex adjoint semisimple G and parabolic P. If a compact complex manifold M has holomorphically trivial tangent bundle, M=\allowbreak{}L/Gamma with Gamma discrete in a complex Lie group L, and every such holomorphic parabolic geometry is the quotient of the source Example2 left invariant Cartan geometry on L. Its Cartan form is h\textasciicircum{}-1dh+\allowbreak{}Ad(h)\textasciicircum{}-1(t ell\textasciicircum{}-1dell), with constant linear injection t whose image complements p. In particular, for a complex torus the geometry is translation invariant. No assertion about non-effective central lifts selected.\par
\noindent\textbf{Difficulty:} H2: A global holomorphic frame alone does not trivialize a parabolic Cartan bundle or force its connection to be invariant. The author constructs a canonical reductive reduction using the holomorphic volume form, root strings and nilradical flows. A holomorphic map to an affine reductive quotient then trivializes that reduction, and compactness forces the remaining Cartan coefficients to be constant in the invariant frame.\par
\noindent\textbf{Editorial boundary note (not author text):} The complete original Theorem3 is selected for effective adjoint rational homogeneous models. The full root/Chevalley setup and root-string calculation, curved structure equations, canonical volume-induced reductive reduction, invariant Example2 and its reconstruction lemma are bundled with the framing/affine-quotient/constant-coefficient proof. Standard reductive affine-quotient and compact parallelizable Lie-group-quotient theorems remain cited prerequisites. The torus specialization belongs to this same main classification; no Calabi-Yau covering, volume-reduction or Weyl helper is separately selected.\par
\noindent\textbf{Source:} \url{https://sigma-journal.com/2011/090/}\quad DOI: \url{https://doi.org/10.3842/SIGMA.2011.090}\par
\noindent\textbf{License and attribution:} Copyright retained by the named authors. Published by SIGMA under CC BY 4.0: \url{https://creativecommons.org/licenses/by/4.0/}. Source: \url{https://sigma-journal.com/about.html}. Adaptation consists of standalone excerpt formatting, cover and numbering restoration; original author language and mathematical text are retained.\par
\noindent\textbf{Editorial symbol notice (not author text):} The effective adjoint model restriction is explicit in the selected scope, so the isotropy representation used in the affine-quotient trivialization is interpreted faithfully. The source leaves exterior-product sign checks routine and gives all substantive reduction/trivialization steps; original equations retained without added proof.\par
\medskip\hrule\medskip\noindent\textbf{Author text begins below.}\par
\setcounter{section}{4}
\setcounter{subsection}{0}
\setcounter{subsubsection}{0}
\setcounter{equation}{0}
\setcounter{section}{3}
\setcounter{subsection}{0}
\setcounter{lemma}{1}
\setcounter{theorem}{1}
\setcounter{example}{1}
% Original source lines 480--868
\section{Rational homogeneous varieties}

Suppose that $G/P$ is a rational homogeneous variety,
so $G$ is a complex semisimple Lie group
and $P$ is a complex parabolic subgroup,
with Lie algebras $\LieG$ and $\LieP$. We can
express $\LieP$ as a sum
of the Cartan subalgebra of $\LieG$ together with
various root spaces, including all of the positive
root spaces. Some negative root spaces will also lie
in $\LieP$. Once we f\/ix the choice
of $\LieG$ and $\LieP$, roots then divide up into
3 categories as follows. The \emph{compact roots} of $\LieG$ are
the roots $\alpha$ of $\LieG$ so that the root spaces of
both $\alpha$ and $-\alpha$ belong to the Lie algebra
of $\LieP$.
All other roots are \emph{noncompact}, and divide
into the \emph{noncompact positive} and \emph{noncompact negative
roots}, according to whether or not their root spaces
lie in $\LieP$.
The Dynkin diagram of $G/P$ is the Dynkin
diagram of $G$ (labelled by simple roots),
with simple roots dotted if they are compact,
and crossed if they are noncompact.

The sum of the root spaces of the noncompact positive roots
is the maximal nilpotent subalgebra, denoted
$\LieN \subset \LieP$.
The sum of the root spaces of the noncompact negative roots
is also a nilpotent subalgebra, denoted
$\LieN^{-} \subset \LieG$, complementary to $\LieP$.
Let $\LieA \subset \LieP$
be the subalgebra spanned by the coroots of the compact roots.
Let $\LieM \subset \LieP$ be the Lie subalgebra generated
by the root space of the compact roots.
The Lie subalgebra
$\LieM \oplus \LieA$ ($\LieN$ resp.)
is the maximal reductive (nilpotent) subalgebra
of $P$; see Knapp \cite{Knapp:2002}.
Let $M$, $A$, $N$ and $N^-$ be the connected subgroups
of $G$ with Lie algebras
$\LieM$, $\LieA$, $\LieN$ and $\LieN^-$.
The groups $M$, $A$, $N$, $N^-$ and $P$
are all algebraic (see Fulton and Harris \cite[p.~382]{Fulton/Harris:1991}).
The splitting
$\LieG = \LieN^{-} \oplus \LieM \oplus \LieA \oplus \LieN$
is $MA$-invariant.

Pick a Chevalley basis $X_{\alpha}$, $H_{\alpha}$
for $\LieG$. Recall (see Serre~\cite{Serre:2001}) that this is a basis
parameterized by roots $\alpha \in \cartan^*$
(with $\cartan \subset \gtot$ a Cartan
subalgebra) for which
\begin{enumerate}\itemsep=0pt
\item $\left[H,X_{\alpha}\right]=\alpha(H) X_{\alpha}$
for each $H \in \cartan$;
\item $\alpha\left(H_{\beta}\right)=
2 \frac{\left<\alpha, \beta\right>}{\left<\beta,\beta\right>}$
(measuring inner products via the Killing form);
\item $\left[H_{\alpha},H_{\beta}\right]=0$;
\item
\[
\left[X_{\alpha},X_{\beta}\right]=
\begin{cases}
H_{\alpha}, & \text{ if } \alpha+\beta=0, \\
N_{\alpha\beta} X_{\alpha+\beta}, & \text{otherwise}
\end{cases}
\]
with
\begin{enumerate}\itemsep=0pt
\item $N_{\alpha\beta}$ an integer,
\item $N_{-\alpha,-\beta}=-N_{\alpha\beta}$,
\item If $\alpha$, $\beta, $ and $\alpha+\beta$ are roots, then
$N_{\alpha\beta}=\pm (p+1)$,
where $p$ is the largest integer for which $\beta-p  \alpha$ is a root,
\item
$N_{\alpha \beta}=0$ if $\alpha+\beta = 0$ or
if any of $\alpha$, $\beta,$ or  $\alpha+\beta$ is not a root.
\end{enumerate}
\end{enumerate}

Consider the 1-forms $\omega^{\alpha}$ dual to the vectors
$X_{\alpha}$ of a Chevalley basis. We use the Killing form to extend $\alpha$ from $\cartan$ to $\gtot$,
by splitting $\gtot = \cartan + \cartan^{\perp}$,
and taking $\alpha=0$ on $\cartan^{\perp}$. The 1-forms~$\omega^{\alpha}$,~$\alpha$
span $\gtot^*$.

Each exterior form in $\Lm{*}{\gtot}^*$ extends uniquely to
a left invariant dif\/ferential form in $\nForms{*}{\Gtot}$, and we will
identify these. These forms determine a basis of left invariant 1-forms
$\omega^{\alpha}$, $\alpha$, and a~basis of left invariant vector
f\/ields $X_{\alpha}$, $H_{\alpha}$. Clearly
\begin{gather*}
d \omega^{\alpha}  = - \alpha \wedge \omega^{\alpha} - \frac{1}{2}
\sum_{\beta + \gamma=\alpha} N_{\beta \gamma} \omega^{\beta} \wedge \omega^{\gamma}, \qquad
d \alpha  =
-
\sum_{\beta} \frac{\left<\alpha,\beta\right>}
{\left<\beta,\beta\right>} \omega^{\beta} \wedge \omega^{-\beta},
\end{gather*}
with sums over all roots.
To be more precise $\omega^{\alpha}$, $\alpha$ is not quite a basis of 1-forms,
since there will be relations among the $\alpha$ 1-forms
in general. To produce a basis, we would have to restrict
to the $\alpha$ 1-forms which are simple roots, but
include all of the $\omega^{\alpha}$ 1-forms,
even for nonsimple~$\alpha$. The basis~$\omega^{\alpha}$,~$\alpha$
is \emph{not} the dual basis to $X_{\alpha}$, $H_{\alpha}$.

\begin{definition}
If $G/P$ is a rational homogeneous variety,
let $\delta = \delta_{G/P}$ be
\[
\delta = \frac{1}{2} \sum_{\alpha}^{\times} \alpha,
\]
where $\sum^{\times}$ means the sum over all noncompact negative roots.
\end{definition}

\begin{lemma}
The Killing form inner product
\(
\left<\delta, \beta\right>
\)
$($where $\delta$ is half the sum of noncompact negative roots, and $\beta$
any root$)$ vanishes
just precisely for $\beta$ a root of the
maximal semisimple subalgebra $\mathfrak{m} \subset \gstruc$.
\end{lemma}
\begin{proof} Knapp \cite[Corollary 5.100, p.~330]{Knapp:2002}
gives a completely elementary proof. We give a~proof along
the same lines, to keep our exposition self-contained.
Pick $\beta$ any positive root. If~$\gamma$ is any
noncompact negative root, and $\left<\gamma,\beta\right> > 0$,
then
\[
\gamma, \gamma-\beta, \gamma-2   \beta, \dots, \gamma-q   \beta = r_{\beta} \gamma
\]
is a string of roots ending in the ref\/lection $r_{\beta}$ of $\gamma$.
To start with, $\gamma$ already contains a positive
multiple of a noncompact negative simple root. Equivalently, $\gamma$
has some negative multiple of a~noncompact positive simple root~$\alpha_1$.
Subtracting the positive
root~$\beta$ can only make the multiple of~$\alpha_1$
larger negative. Therefore the entire string consists
of noncompact negative roots.

If we have an entire $\beta$-string of noncompact
negative roots, for a positive
root $\beta$, clearly
\begin{gather*}
\left<r_{\beta} \gamma,\beta\right>
 = -\left<r_{\beta} \gamma, r_{\beta} \beta\right>
 =
-\left<\gamma,\beta\right>.
\end{gather*}
Therefore $\left<\gamma,\beta\right>$ cancels with $\left<r_{\beta}
\gamma,\beta\right>$ in the sum $\left<\delta,\beta\right>$. Hence
the entire string cancels out of that sum.

It follows that
\begin{equation}\label{eqn:Sum}
\left<\delta,\beta\right>=
\sum_{\gamma} \left<\gamma,\beta\right>
\end{equation}
where the sum is over noncompact negative roots $\gamma$
for which both $\left<\gamma,\beta\right> \le 0$
and for which the other end of the $\beta$-string
through $\gamma$ is noncompact positive.
Of course, the terms with  $\left<\gamma,\beta\right> = 0$
cancel out too, so the sum~(\ref{eqn:Sum}) is over
noncompact negative roots $\gamma$ for
which both $\left<\gamma,\beta\right> < 0$
and for which the other end of the $\beta$-string
through~$\gamma$ is noncompact positive. In particular,
the sum (\ref{eqn:Sum}) is a sum of negative terms.
But there might not be any terms.

If $\beta$ is a compact root, then clearly ref\/lection in
$\beta$ preserves the roots belonging to the parabolic subalgebra
$\gstruc$, and therefore preserves the noncompact negative roots.
So the noncompact negative roots will all lie in $\beta$-strings, and
$\left<\delta,\beta\right>=0$ for these roots. On the other hand, if~$\beta$ is a noncompact root, then $\beta$
is either positive or negative. We can assume that~$\beta$ is positive, since we only need to show that $\left<\delta,\beta\right> \ne 0$. Take
$\gamma=-\beta$, to see that the sum~(\ref{eqn:Sum}) has at least
one negative term.
\end{proof}

\section{Parabolic geometries}

The standard reference on parabolic
geometries is \v{C}ap and Slovak \cite{Cap/Slovak:2009};
we will use the standard def\/initions, as in their
book, which are far too long to put into this paper.
Suppose that $E \to M$ is a holomorphic
parabolic geometry, with some model $G/P$.
Very similar structure equations hold for any
holomorphic parabolic geometry with the same model. Indeed,
the Cartan connection is a 1-form valued in the Lie
algebra $\mathfrak{g}$ of $G$, so splits into a sum
of 1-forms~$\omega^{\alpha}$ and~$\alpha$ from the
decomposition of $\mathfrak{g}$ into root spaces.
From the def\/inition of a Cartan geometry, the
Cartan connection satisf\/ies the same structure
equations as the Maurer--Cartan form on the model,
but with semibasic curvature correction terms, so
\begin{gather*}
d \omega^{\alpha}  = - \alpha \wedge \omega^{\alpha} - \frac{1}{2}
\sum_{\beta + \gamma=\alpha} N_{\beta \gamma} \omega^{\beta} \wedge \omega^{\gamma}
+
\sum^{\times}_{\beta, \gamma}
\kappa^{\alpha}_{\beta \gamma} \omega^{\beta} \wedge \omega^{\gamma},
\\
d \alpha  =
-
\sum_{\beta} \frac{\left<\alpha,\beta\right>}
{\left<\beta,\beta\right>} \omega^{\beta} \wedge \omega^{-\beta}
+
\sum^{\times}_{\beta, \gamma}
\lambda^{\alpha}_{\beta \gamma} \omega^{\beta} \wedge \omega^{\gamma},
\end{gather*}
where the $\kappa$ and $\lambda$ terms are Cartan geometry
curvature terms, so they vanish except possibly for
$\beta$ and $\gamma$ noncompact negative roots, and
once again $\displaystyle{\sum^{\times}}$ means
the sum over noncompact negative roots.

It is vital in the following that, even if we
work on a manifold where we have imposed some
relations on these 1-forms, we will still use the
Killing form on the original Lie algebra $\gtot$
to compute inner products $\left<\alpha,\beta\right>$.
This is our only notational ambiguity.

\section{Proofs of the theorems}

Replacing our Calabi--Yau manifold by a f\/inite covering space
if needed, we can assume that it
bears a nowhere-vanishing holomorphic volume
form. We then derive our theorem from the following stronger theorem:

\begin{theorem}\label{theorem:Reduction}
If a complex manifold bears a holomorphic parabolic geometry
and a holomorphic volume form, then it admits a canonical
holomorphic reduction of the structure group of the parabolic
geometry to a reductive algebraic group.
\end{theorem}
\begin{proof}
Suppose that $E \to M$ is a holomorphic parabolic
geometry modelled on $\Gtot/\Gstruc$,
and $\sigma$ a~holomorphic volume form on $M$.
Pick a Chevalley basis. Let
\begin{gather*}
\Omega  = \bigwedge^{\times}_{\alpha} \omega^{\alpha}, \qquad
\delta  = \frac{1}{2} \sum^{\times}_{\alpha} \alpha,
\end{gather*}
where the wedge product and sum are over noncompact negative roots.
The sign of $\Omega$ depends on a choice of ordering
of the noncompact negative roots, but any ordering can be chosen, as
long as we are consistent.

We claim that $d \Omega = - 2 \delta \wedge \Omega$.
Order the noncompact negative
roots arbitrarily as $\alpha_1, \alpha_2, \dots$.
Expand out $d \Omega$:
\[
d \Omega
=
 \sum_{j}
(-1)^{j+1}
\bigwedge_{i < j} \omega^{\alpha_i}
\wedge d \omega^{\alpha_j}
\wedge
\bigwedge_{i > j} \omega^{\alpha_i},
\]
by passing the exterior derivative operator along
the various factors, hitting one $\omega^{\alpha}$
at a time, and sticking a suitable $\pm$ sign in front.
Plug in the equation for $d \omega^{\alpha}$,
and the curvature terms all vanish because they
occur in pairs of $\omega^{\beta} \wedge \omega^{\gamma}$,
and at least one of these $\omega^{\beta}$ or $\omega^{\gamma}$
is still present in factors in that term.
If we f\/ind a term like $N_{\beta \gamma} \omega^{\beta}
\wedge \omega^{\gamma}$ in $d \omega^{\alpha}$,
then we must have $\alpha=\beta+\gamma$.
Since $\alpha$ is a noncompact negative root, at least one
of $\beta$ and $\gamma$ must be as well. Neither can be
equal to $\alpha$, since $N_{0\alpha}=0$. Therefore
each such term drops out of $d \Omega$. The reader
now only has to check the signs to see that
$d \Omega = -2 \delta \wedge \Omega$.

Our nonzero section $\sigma$ of the canonical
bundle of $M$ can be pulled back to $E$ as
\(
\sigma = s \Omega,
\)
for a unique nowhere-vanishing function $s : E \to \C{}$.
If $\sigma$ is holomorphic, then
\begin{gather*}
0  =
d\sigma
 =
ds \wedge \Omega + s   d \Omega
=
\left(ds - 2 s \delta \right) \wedge \Omega.
\end{gather*}
Let $P_0 \subset P$ be the subgroup of $P$
acting trivially on
$\Lambda^{\text{top},0}\left(\mathfrak{g}/\mathfrak{p}\right)$.
Let $E' \subset E$ be the set of points at which
$s=1$. Then $E' \subset E$ is a smooth hypersurface since $ds \ne 0$
on tangent spaces of~$E$ along~$E'.$ Clearly
$E'$ is a principal right $P_0$-bundle.

On $E'$, $\delta \wedge \Omega = 0$.
Therefore $\delta$ is semibasic on $E'$:
\[
\delta = \sum^{\times}_{\alpha} t_{\alpha} \omega^{\alpha},
\]
a sum over noncompact negative roots $\alpha$, for some
functions $t_{\alpha} : E' \to \mathbb{C}.$
Taking exterior derivative, we f\/ind
\begin{gather*}
0 =
d \left(\delta - \sum _{\alpha} t_{\alpha} \omega^{\alpha} \right)
 = \sum_{\alpha}
\left(
d \alpha
-
dt_{\alpha} \wedge \omega^{\alpha}
- t_{\alpha} d \omega^{\alpha}
\right)
\\
\phantom{0}{} =
-
\sum_{\beta} \frac{\left<\delta,\beta\right>}
{\left<\beta,\beta\right>} \omega^{\beta} \wedge \omega^{-\beta}
-
\sum_{\alpha}
\left(dt_{\alpha} - t_{\alpha} \alpha\right) \wedge \omega^{\alpha}
\\
\phantom{0=}{} - \frac{1}{2}
\sum_{\alpha} t_{\alpha}
\sum_{\beta + \gamma=\alpha} N_{\beta \gamma} \omega^{\beta} \wedge \omega^{\gamma}
\pmod{\text{semibasic terms}}.
\end{gather*}
In particular, for any noncompact negative root $\alpha$,
\[
\LieDer_{X_{-\alpha}} t_{\alpha} = 2 \frac{\left<\delta,\alpha\right>}{\left<\alpha,\alpha\right>}.
\]
Since $-\alpha$ is a positive root, the corresponding root vector
lies in $\mathfrak{p}$. Moreover, this root vector lies in the
nilpotent radical of $\mathfrak{p}$, since it is a positive root.
The nilpotent radical acts trivially on~$\Lambda^{\text{top},0}\left(\mathfrak{g}/\mathfrak{p}\right)$,
as nilpotent groups have no nontrivial 1-dimensional representations.
Every vector in $\mathfrak{p}$ gives rise to a vector f\/ield
giving the associated inf\/initesimal action, and for the root
vector of $-\alpha$, this vector f\/ield is $X_{-\alpha}$,
by def\/inition of a Cartan geometry. Since the root
vector lies in the Lie algebra of $P_0$, the vector f\/ield
$X_{-\alpha}$ generates a 1-parameter subgroup of $P_0$.
The vector f\/ield $X_{-\alpha}$ is therefore
tangent to the f\/ibers of $E' \to M$.

On the f\/ibers the vector f\/ield $X_{-\alpha}$ is a left invariant vector f\/ield.
Therefore $X_{-\alpha}$ is complete. Starting at any point of $E'$,
we can move in the direction $X_{-\alpha}$ of the nilpotent part
of the structure group, altering the value of $t_{\alpha}$ at a constant rate
until it reaches $0$. Indeed $t_{\alpha}$ is acted on by the nilradical
of the structure group as translations in the left action on $E'$. The set of points
$E'' \subset E'$ on which all $t_{\alpha}$ vanish is a smooth embedded
submanifold, because its tangent space is cut out by equations
\[
\frac{\left<\delta,\alpha\right>}{\left<\alpha,\alpha\right>} \omega^{-\alpha} = \text{semibasic},
\]
for all noncompact negative roots $\alpha.$ The structure group is
reduced to
a reductive algebraic group, since we have eliminated the nilradical
of the original structure group, leaving only the root spaces
$\alpha$ for which neither $\alpha$ nor $-\alpha$ is
noncompact negative, i.e.
the root spaces of the maximal reductive subgroup $MA$
of the structure group $P$. We have also eliminated the
part of~$A$ which acts nontrivially on the holomorphic
volume forms, so our structure group is now~$MA^0$,
with~$A^0$ the subgroup of~$A$ f\/ixing a volume form on
$\mathfrak{g}/\mathfrak{p}$.
\end{proof}
\setcounter{section}{7}
\setcounter{subsection}{0}
\setcounter{subsubsection}{0}
\setcounter{equation}{1}
\setcounter{section}{6}
\setcounter{subsection}{0}
\setcounter{example}{1}
% Original source lines 899--1049
\section{Parabolic geometries on tori}

\begin{example}\label{example:TranslationInvariant}
Suppose that $L$ is a Lie group with Lie algebra $\LieL$,
and write the left invariant Maurer--Cartan 1-form
on $L$ as $\ell^{-1}   d\ell$.
Suppose that $G$ is a Lie group and $H \subset G$
is a closed subgroup. Take any linear injection $t : \LieL \to \LieG$
so that the image is complementary to $\LieH$.
Then let $M=L$ and $E=M \times H$. Let $\omega \in \nForms{1}{E} \otimes  \LieG$
be the 1-form
\[
\omega = h^{-1}   dh + \Ad(h)^{-1}\big(t   \ell^{-1}   d\ell\big).
\]
It is easy to check that $\omega$ is the Cartan
connection of Cartan geometry on $L$ modelled on $G/H$.
The group $L$ acts as Cartan geometry automorphisms
by the obvious action. The curvature is
\[
d \omega + \frac{1}{2}\left[\omega,\omega\right]
=
\frac{1}{2}
\Ad(h)^{-1}\left(
\left[
t \ell^{-1}   d\ell,
t \ell^{-1}   d\ell
\right]
-
t
\left[
\ell^{-1}   d\ell,
\ell^{-1}   d\ell
\right]
\right).
\]
In particular, the Cartan geometry is f\/lat if and only if
$t$ is a Lie algebra homomorphism. See~\cite{Hammerl:2007}
for details. Clearly the group $\Aut(L) \times \Aut(G,H)$
acts as isomorphisms of these geometries.
\end{example}


\begin{lemma} Every left invariant Cartan geometry on a
Lie group is equivariantly isomorphic to one
constructed as in Example~{\rm \ref{example:TranslationInvariant}}.
\end{lemma}
\begin{proof}
Take any Cartan geometry $E \to L$ on a Lie group $L$.
Suppose that $L$ acts on $E$ lifting its left action
on itself, commuting with the $H$-action,
and preserving the Cartan connection.
Pick any point $e_0 \in E$. Map
\[
\left(\ell,h\right) \in L \times H \mapsto \ell e_0 h.
\]
Clearly this is an isomorphism of principal bundles,
so from now on we take $E=L \times H$.
Unwinding the def\/inition of a Cartan connection immediately yields
that every Cartan connection on the bundle $L \times H \to L$ has the form
\[
\omega = h^{-1}   dh + \Ad(h)^{-1} \left(\gamma\right),
\]
where $\gamma$ is a 1-form on $L$ valued in $\LieG$,
so that $\gamma+\LieH$ is a linear isomorphism,
i.e.\ $\gamma_{\ell} : T_{\ell} L \to \gtot$ is a
linear injection, for each $\ell \in L$, and
$T_{\ell} L \to \gtot \to \gtot/\gstruc$ is a linear isomorphism.
By translation invariance
\[
\gamma = t \ell^{-1}   d\ell,
\]
for a unique linear map $t$ and
$t \to \gtot \to \gtot/\gstruc$ is a linear isomorphism.
\end{proof}

\begin{theorem}\label{theorem:translation}
Every holomorphic parabolic geometry on any complex torus is translation
inva\-riant, and obtained by the construction of
Example~{\rm \ref{example:TranslationInvariant}}
$($where the Lie group $L$ is the complex torus itself$)$.
More generally, if $M$ is any compact complex
manifold with holomorphically trivial tangent bundle,
then $M=L/\Gamma$ for some discrete
subgroup $\Gamma \subset L$ of a complex
Lie group~$L$. Every holomorphic parabolic geometry
on $M$ is obtained by the construction of
Example~{\rm \ref{example:TranslationInvariant}}
applied to~$L$, and then quotiented by~$\Gamma$.
\end{theorem}

\begin{proof}
If $M$ is a compact complex manifold with holomorphically
trivial tangent bundle, then any basis of the
holomorphic tangent space $T_m M$ at any point $m \in M$
extends to a framing by holomorphic vector f\/ields, say $X_1, X_2, \dots, X_n$.
These then generate a transitive complex Lie group action,
say of a complex Lie group $L$. Brackets of these
holomorphic vector f\/ields can be rewritten in terms
of the vector f\/ields themselves
\[
\left[X_i,X_j\right]=\sum c^k_{ij} X_k,
\]
since the $X_k$ form a basis. The holomorphic functions
$c^k_{ij}$ are constant because $M$ is compact. Therefore $L$ has the same
dimension as $M$, and so $M = L/\Gamma$ \cite{Wang:1954}.

Following the proof of Theorem~\ref{theorem:Reduction}, the structure group of any
holomorphic parabolic geometry $E \to M$ reduces to a reductive group, $G''$
on some subbundle $E'' \subset E$. The Cartan connection~$\omega$ splits into a sum corresponding to the splitting
of $\mathfrak{g}$ into $G''$-invariant subspaces,
and~$\omega''$ (the part valued in $\mathfrak{g}''$)
is a connection form for $E'' \to M$. Take a global
holomorphic coframing on~$M$, i.e.\ a set of linearly independent 1-forms
$\xi^{\alpha}$ forming a basis of each cotangent space of~$M$, for
$\alpha$ varying over noncompact negative roots.
Def\/ine a map $e \in E'' \to h
\in \GL{\gtot/\gstruc}$, by $\omega^{\alpha}=h^{\alpha}_{\beta}
\xi^{\beta}$ (for $\alpha$ and $\beta$ varying over
noncompact negative roots),
and $h(e) =  (h^{\alpha}_{\beta} )$ in the basis~$X_{\alpha}$ for the sum of noncompact negative
root spaces. Under right
$G''$-action,
\[
h\left(r_{g}e\right)=g^{-1}h(e)
\]
for $g \in G''.$ Therefore the quotient map $E \!\to\! \GL{\gtot/\gstruc}/G''$
descends to a map $M \!\to\! \GL{\gtot/\gstruc}/G''$.
The quotient $\GL{\gtot/\gstruc}/G''$ is an af\/f\/ine variety:
see Mumford et al.~\cite[Theorem~1.1, p.~27]{Mumford/Fogarty/Kirwan:1994}
and Procesi~\cite[Theorem~2, p.~556]{Procesi:2007}.
Af\/f\/ine coordinate functions will pull back to functions on $M$,
and therefore must be constant.
Therefore the map $M \to \GL{\gtot/\gstruc}/G''$ is constant.
We have an isomorphism
\[
e \in E'' \to \left(\pi(e),h(e)\right) \in M \times G'',
\]
trivializing the bundle $E''$. We can therefore assume
that $E'' = M \times G''$ and $E = M \times H$.
Again unwinding the def\/inition of a Cartan
geometry,
\[
\omega = h^{-1}   dh + \Ad(h)^{-1}\gamma,
\]
and $\gamma \in \nForms{1}{M} \otimes \LieG$.
The functions $X_i \hook \gamma$
are holomorphic functions on $M$,
so constant, so $\gamma$ pulls back to $L$ to be
$\gamma=t \ell^{-1}   d\ell$ for some constant
linear map $t : \LieL \to \LieG$.
\end{proof}
\begin{thebibliography}{99}
\bibitem[7]{Cap/Slovak:2009}
\v{C}ap A., Slov{\'a}k J.,
Parabolic geometries.~I.~Background and general theory,
{\em Mathematical Surveys and Monographs}, Vol.~154,
American Mathematical Society, Providence, RI, 2009.


\bibitem[11]{Fulton/Harris:1991}
 Fulton W.,  Harris J.,
  Representation theory. A~f\/irst course,
  {\em Graduate Texts in   Mathematics}, Vol.~129,
 Springer-Verlag, New York, 1991.

\bibitem[13]{Hammerl:2007}
Hammerl M.,
Homogeneous Cartan geometries,
{\em Arch. Math. (Brno)} {\bf 43} (2007), 431--442,
\href{http://arxiv.org/abs/math.DG/0703627}{math.DG/0703627}.

\bibitem[20]{Knapp:2002}
 Knapp A.W.,
Lie groups beyond an introduction, 2nd ed.,
 {\em Progress   in Mathematics}, Vol.~140,
  Birkh\"auser Boston Inc., Boston, MA, 2002.

\bibitem[29]{Mumford/Fogarty/Kirwan:1994}
Mumford D., Fogarty J.,  Kirwan F.,
 Geometric invariant theory,  3rd ed., {\em Ergebnisse der
  Mathematik und ihrer Grenzgebiete (2)}, Vol.~34,
 Springer-Verlag, Berlin,   1994.

\bibitem[30]{Procesi:2007}
 Procesi C.,
Lie groups. An approach through invariants and representations, {\it Universitext}, Springer, New York, 2007.

\bibitem[31]{Serre:2001}
 Serre J.-P.,
 Complex semisimple Lie algebras,
{\it Springer Monographs in Mathematics}, Springer-Verlag, Berlin, 2001.


\bibitem[33]{Wang:1954}
 Wang H.-C.,
 Closed manifolds with homogeneous complex structure,
\href{http://dx.doi.org/10.2307/2372397}{{\em Amer.~J. Math.}} {\bf 76} (1954), 1--32.

\end{thebibliography}
\end{document}
